%generare il pdf con il comando: pdflatex main.tex
\documentclass[a4paper, oneside, openany]{article}
\usepackage{../template/sos}
\usepackage[dvipsnames]{xcolor}
%%%%%%%%%%%%%%%%%%%%%%%%%%%%%%%%%%%%%%%%%%%%%%%%%%%%%%%%%%%%%%%%%%%%%%%%%%%%%%%%%%%%%%%%%%%%%%%%%%%%%%%
%MACRO
%%%%%%%%%%%%%%%%%%%%%%%%%%%%%%%%%%%%%%%%%%%%%%%%%%%%%%%%%%%%%%%%%%%%%%%%%%%%%%%%%%%%%%%%%%%%%%%%%%%%%%%

\newcommand{\Titolo}{Verbale Interno}

\newcommand{\Gruppo}{Gruppo Sons Of SWE}

\newcommand{\Redazione}{Thiella Eleonora}

\newcommand{\ACapoRedazione}{Thiella Eleonora}

\newcommand{\Verifica}{Cavallin Giovanni}

\newcommand{\Approvazione}{Panozzo Stefano}

\newcommand{\Distribuzione}{Vardanega Tullio \newline Cardin Riccardo \newline Gruppo Sons Of SWE}

\newcommand{\Uso}{Interno}

\newcommand{\Data}{3 Marzo 2018}

\newcommand{\NomeProgetto}{Progetto Marvin}

\newcommand{\Mail}{sonsofswe.swe@gmail.com}

\newcommand{\DescrizioneDoc}{Questo documento contiene il secondo \VI\ped{G} svoltosi in data 8 marzo 2018 dal gruppo Gruppo Sons Of SWE riguardo al capitolato Marvin.}
%%%%%%%%%%%%%%%%%%%%%%%%%%%%%%%%%%%%%%%%%%%%%%%%%%%%%%%%%%%%%%%%%%%%%%%%%%%%%%%%%%%%%%%%%%%%%%%%%%%%%%%
%FINE MACRO
%%%%%%%%%%%%%%%%%%%%%%%%%%%%%%%%%%%%%%%%%%%%%%%%%%%%%%%%%%%%%%%%%%%%%%%%%%%%%%%%%%%%%%%%%%%%%%%%%%%%%%%
%INIZIO INCLUDE PACCHETTI
%%%%%%%%%%%%%%%%%%%%%%%%%%%%%%%%%%%%%%%%%%%%%%%%%%%%%%%%%%%%%%%%%%%%%%%%%%%%%%%%%%%%%%%%%%%%%%%%%%%%%%%
\usepackage{graphbox}
% permette di modificare i margini
\usepackage[top=3.1cm, bottom=3.1cm, left=2.2cm, right=2.2cm]{geometry}

\usepackage{lastpage} %info sul # dell'ultima pagina del documento
\usepackage{fancyhdr} %per modificare dimensioni,margini, intestazioni e righe a piè di pagina
\fancypagestyle{plain}{
  % cancella tutti i campi di intestazione e piè di pagina
  \fancyhf{}

  \lhead{ %in alto a sx
  % da modificare quando avremo un logo
  	Metterci un logo quando ne avremo uno
  }

  \rhead{
  	\Mail{}
  }


  \lfoot{ %piè di pagina a sx
    \Titolo{} \ - \textit{\Gruppo{}}
  }
  \rfoot{Pagina \thepage{} di \pageref{LastPage}} %es: pag: 4 di 10

  %linea orizzontale alle posizioni top e bottom della pagina (se è 0, non c'è la linea)
  \renewcommand{\headrulewidth}{0.3pt}  
  \renewcommand{\footrulewidth}{0.3pt}
}
\pagestyle{plain}

%\usepackage{calc} %introduce la notazione infissa per le op. aritmetiche interne a LaTeX

\usepackage[utf8]{inputenc}
\usepackage[T1]{fontenc}
\usepackage[italian]{babel} %il documento è in italiano
%\usepackage{textcomp} %The pack­age sup­ports the Text Com­pan­ion fonts, which pro­vide many text sym­bols
%(such as baht, bul­let, copy­right, mu­si­cal­note, onequar­ter, sec­tion, and yen), in the TS1 en­cod­ing.

\usepackage{graphicx}       %permette di inserire delle immagini
\usepackage{caption}        %numerazione figure e loro descrizione testuale
\usepackage{subcaption}     %sottofigure numerabili
\usepackage{float}  %permette di inserire un # qualsiasi di figure fluttuanti

\usepackage{rotating} %permette di ruotare le immagini
%\usepackage{changepage} %utile se c'è bisogno di aggiustare margini per centrare figure

%package utili per la math mode ( $ ... $ o \[ ... \] )
\usepackage{amsmath}
\usepackage{amssymb}
\usepackage{amsfonts}
%\usepackage{euler}    %font 'ams euler', lo stesso di 'Concrete Mathematics' di Knuth
\usepackage{amsthm}
\usepackage{mathtools}

\usepackage{listings} %permette di inserire degli spezzoni di codice

\usepackage{tikz} %disegno di immagini vettoriali a schermo. Utile per grafi
\usetikzlibrary{arrows.meta}
\usetikzlibrary{graphs}
\usetikzlibrary{arrows}
%\usepackage{tikz-uml} %serve per disgnare l'UML, fantastica guida:
%https://perso.ensta-paristech.fr/~kielbasi/tikzuml/var/files/doc/tikzumlmanual.pdf
%download package: http://perso.ensta-paristech.fr/~kielbasi/tikzuml/

%package per le tabelle
\usepackage{booktabs} %permette di poter usare delle liste nelle tabelle
\usepackage{tabularx} 
\usepackage{longtable} %una tabella può continuare su più pagine
\usepackage{multirow} %utile per visualizzare una cella su più righe
%\usepackage{multicolumn} %cella su più colonne
%\usepackage[table]{xcolor} %rende disponibile l'utilizzo di un colore per lo sfondo
                        %delle celle di una tabella

%crea una cella per le tabelle in grado di andare a capo con \newline
%https://tex.stackexchange.com/questions/12703/how-to-create-fixed-width-table-columns-with-text-raggedright-centered-raggedlef
\usepackage{array}
\newcolumntype{L}[1]{>{\raggedright\let\newline\\\arraybackslash\hspace{0pt}}m{#1}}
\newcolumntype{C}[1]{>{\centering\let\newline\\\arraybackslash\hspace{0pt}}m{#1}}
\newcolumntype{R}[1]{>{\raggedleft\let\newline\\\arraybackslash\hspace{0pt}}m{#1}}


%indice con i puntini
\usepackage{tocloft}
\renewcommand\cftsecleader{\cftdotfill{\cftdotsep}}

%http://ctan.mirror.garr.it/mirrors/CTAN/macros/latex/contrib/appendix/appendix.pdf
\usepackage{appendix} %aggiunge dei comandi per l'appendice
\usepackage{parskip} %aiuta LaTeX a trovare il miglior stile per i page break
\setcounter{secnumdepth}{5} % numera i sottoparagrafi
\setcounter{tocdepth}{5} %aggiunge all'indice i sottoparagrafi
%\usepackage{titlesec} %\begin{paragraph} si può usare come subsubsubsection!


\usepackage{breakurl}%\url{...} può continare alla linea successiva. (si può andare a capo)


\usepackage[colorlinks=true]{hyperref}
\hypersetup{
    colorlinks=true,
    citecolor=black,
    filecolor=black,
    linkcolor=black, % colore dei link interni
    urlcolor=Maroon  % colore dei link interniesterni
}

%impostazioni per il codice che deve finire dentro a
%\begin{lstlisting}

\definecolor{listinggray}{gray}{0.9}
\definecolor{lbcolor}{rgb}{0.9,0.9,0.9}
\lstset{
backgroundcolor=\color{lbcolor},
    tabsize=4,    
%   rulecolor=,
    language=[GNU]C++,
    basicstyle=\scriptsize,
    upquote=true,
    aboveskip={1.5\baselineskip},
    columns=fixed,
    showstringspaces=false,
    extendedchars=true,
    inputencoding=utf8,
    breaklines=true,
    prebreak = \raisebox{0ex}[0ex][0ex]{\ensuremath{\hookleftarrow}},
    frame=single,
    numbers=left,
    showtabs=false,
    showspaces=false,
    showstringspaces=false,
    identifierstyle=\ttfamily,
    keywordstyle=\color[rgb]{0,0,1},
    commentstyle=\color[rgb]{0.026,0.112,0.095},
    stringstyle=\color[rgb]{0.627,0.126,0.941},
    numberstyle=\color[rgb]{0.205, 0.142, 0.73},
%        \lstdefinestyle{C++}{language=C++,style=numbers}’.
}
\lstset{
  backgroundcolor=\color{lbcolor},
  tabsize=4,
  language=C++,
  captionpos=b,
  tabsize=3,
  frame=lines,
  numbers=left,
  numberstyle=\tiny,
  numbersep=5pt,
  breaklines=true,
  showstringspaces=false,
  basicstyle=\footnotesize,
  identifierstyle=\color{magenta},
  keywordstyle=\color[rgb]{0,0,1},
  commentstyle=\color{orange},
  stringstyle=\color{red}
}
%%%%%%%%%%%%%%%%%%%%%%%%%%%%%%%%%%%%%%%%%%%%%%%%%%%%%%%%%%%%%%%%%%%%%%%%%%%%%%%%%%%%%%%%%%%%%%%%%%%%%%%
%FINE INCLUDE
%%%%%%%%%%%%%%%%%%%%%%%%%%%%%%%%%%%%%%%%%%%%%%%%%%%%%%%%%%%%%%%%%%%%%%%%%%%%%%%%%%%%%%%%%%%%%%%%%%%%%%%
%INIZIO DOCUMENTO
%%%%%%%%%%%%%%%%%%%%%%%%%%%%%%%%%%%%%%%%%%%%%%%%%%%%%%%%%%%%%%%%%%%%%%%%%%%%%%%%%%%%%%%%%%%%%%%%%%%%%%%
\begin{document}
%%%%%%%%%%%%%%%%%%%%%%%%%%%%%%%%%%%%%%%%%%%%%%%%%%%%%%%%%%%%%%%%%%%%%%%%%%%%%%%%%%%%%%%%%%%%%%%%%%%%%%%
%INIZIO COPERTINA
%%%%%%%%%%%%%%%%%%%%%%%%%%%%%%%%%%%%%%%%%%%%%%%%%%%%%%%%%%%%%%%%%%%%%%%%%%%%%%%%%%%%%%%%%%%%%%%%%%%%%%%
\begin{titlepage}
  \begin{center}

  \begin{center}
    %sostituire con immagine copertina del gruppo
   % \centerline{\includegraphics[scale=0.24]{img/logo.png}}
  \end{center}
  
  \vspace{1cm}

  \begin{Huge}
    \textbf{\Titolo{}} \\
  \end{Huge}

  \vspace{9pt}  
  
  \begin{large}
  	\textbf{\Gruppo{}}\ - \textbf{\NomeProgetto{}}\\%\Data{}
  	\vspace{6pt}
  	\Mail{}
  \end{large}	  
  
  \vspace{15pt}

  \bgroup
  \def\arraystretch{1.3}
   \centering
   \begin{tabular}{c|L{5cm}}
      \multicolumn{2}{c}{\textbf{Informazioni sul documento} } \\ \hline
      %Versione &  1.0.0\\
      \textbf{Redazione} & \ACapoRedazione{} \\
      \textbf{Verifica} & \Verifica{} \\ 
      \textbf{Approvazione} & \Approvazione{} \\
      \textbf{Uso} & \Uso \\
      \textbf{Distribuzione} & \Distribuzione{}
    \end{tabular}
  \egroup

  \vspace{15pt}

  \begin{center}
    \textbf{Descrizione\\}
    \DescrizioneDoc{}
  \end{center}

  \end{center}
  \end{titlepage}
%%%%%%%%%%%%%%%%%%%%%%%%%%%%%%%%%%%%%%%%%%%%%%%%%%%%%%%%%%%%%%%%%%%%%%%%%%%%%%%%%%%%%%%%%%%%%%%%%%%%%%%
%FINE COPERTINA
%%%%%%%%%%%%%%%%%%%%%%%%%%%%%%%%%%%%%%%%%%%%%%%%%%%%%%%%%%%%%%%%%%%%%%%%%%%%%%%%%%%%%%%%%%%%%%%%%%%%%%%
%SOMMARIO
%%%%%%%%%%%%%%%%%%%%%%%%%%%%%%%%%%%%%%%%%%%%%%%%%%%%%%%%%%%%%%%%%%%%%%%%%%%%%%%%%%%%%%%%%%%%%%%%%%%%%%%
%PARAGRAFI
%%%%%%%%%%%%%%%%%%%%%%%%%%%%%%%%%%%%%%%%%%%%%%%%%%%%%%%%%%%%%%%%%%%%%%%%%%%%%%%%%%%%%%%%%%%%%%%%%%%%%%%
\definecolor{bluSOS}{RGB}{48, 84, 150}
\definecolor{greySOS}{RGB}{209, 222, 223}
\section*{Registro delle modifiche}
{
	\rowcolors{2}{greySOS}{white}
	\renewcommand{\arraystretch}{1}
	\centering
	\begin{longtable}{| c| c | C{4cm} | c | c |}
		\hline
		\rowcolor{bluSOS}
		\textcolor{white}{\textbf{Versione}} & \textcolor{white}{\textbf{Data}} & \textcolor{white}{\textbf{Descrizione}} & \textcolor{white}{\textbf{Autore}} & \textcolor{white}{\textbf{Ruolo}}\\
		\hline
		1.0.0 & 2018-04--- & Approvazione & Stefano Panozzo  & Reponsabile \\
		\hline
		0.1.0 & 2018-03-10 & Verificazione & Giovanni Cavallin  & Verificatore \\
		\hline
		0.0.1 & 2018-03-07 & Stesura della sezione \emph{Informazioni generali} e \emph{Riassunto della riunione} & Eleonora Thiella & Analista \\
		
	\end{longtable}

}

%\emph{Analisi dei Rischi}, \emph{Preventivo} ed \emph{Attualizzazione dei rischi}. Create le appendici \emph{A Attuazione dei rischi} e \emph{B organigramma}

\newpage
\tableofcontents
\newpage
\section{Informazioni generali}
	\begin{itemize}
		\item \textbf{Data riunione}: 2018-03-08;
		\item \textbf{Ora inizio riunione}: 10:30;
		\item \textbf{Ora fine riunione}: 12:30;
		\item \textbf{Durata riunione}: 2 Ore;
		\item \textbf{Luogo di incontro}: Luf1, Plesso Paolotti;
		\item \textbf{Oggetto di discussione}: Approfondimento dei temi principali trattati nel \emph{capitolato}\ped{G} scelto, dei documenti da consegnare per la \textcolor{red}{Revisione dei Requisiti} e i ruoli che ogni membro del gruppo deve ricoprire;
		\item \textbf{Moderatore}: Stefano Panozzo;
		\item \textbf{Segretario}: Eleonora Thiella;
		\item \textbf{Partecipanti}: Giovanni Cavallin, Eleonora Thiella, Giovanni Dalla Riva, Stefano Panozzo, Andrea Favero e Federico Caldart;
	\end{itemize}

\section{Riassunto della riunione}
	\subsection{Descrizione}
	Prima della riunione ogni membro si è informato in modo differente sui temi di \emph{blockchain}\ped{G} ed \emph{Ethereum}\ped{G} e all'incontro ha fatto una breve presentazione al resto del gruppo, in modo che ognuno acquisisse una maggior conoscenza di questi argomenti.
	\\Il \emph{team}\ped{G} si è poi informato sui documenti da presentare alla \textcolor{red}{Revisione dei Requisiti}, nello specifico di cosa essi trattano, le informazioni che devono contenere e come devono essere redatti.
	\\Per quanto riguarda la suddivisione del carico di lavoro, il gruppo ha preso in considerazione gli impegni personali di tutti i membri e la familiarità di ognuno di essi con le tecnologie da utilizzare.
	\subsection{Decisioni prese}
		\begin{itemize}
			\item \textbf{VI2.1}: ad ogni membro del gruppo è stato assegnato un ruolo per la stesura dei documenti:
			\begin{itemize}
				\item Stefano Panozzo: \RdP{}\ped{G};
				\item Giovanni Dalla Riva: \ana{}\ped{G};
				\item Eleonora Thiella: \ana{};
				\item Federico Caldart: \ana{};
				\item Lorenzo Menegon: \ver{}\ped{G};
				\item Giovanni Cavallin: \amm{}\ped{G};
				\item Andrea Favero: \amm{}.
			\end{itemize}
		\end{itemize}


\end{document}